\begin{figure}[t]
\centering
% Shared color TikZ styles for the paper figures.
\colorlet{paperBlue}{blue!9!white}
\colorlet{paperBlueLine}{blue!55!black}
\colorlet{paperTeal}{teal!10!white}
\colorlet{paperTealLine}{teal!55!black}
\colorlet{paperAmber}{orange!15!white}
\colorlet{paperAmberLine}{orange!60!black}
\colorlet{paperRose}{red!9!white}
\colorlet{paperRoseLine}{red!55!black}
\colorlet{paperViolet}{violet!10!white}
\colorlet{paperVioletLine}{violet!55!black}
\colorlet{paperInk}{black!72}

\tikzset{
  paperfig/.style={
    font=\small,
    >=Stealth,
    line cap=round,
    line join=round
  },
  paperline/.style={
    -{Stealth[length=2.1mm,width=1.6mm]},
    line width=0.58pt,
    draw=paperInk
  },
  paperdash/.style={
    paperline,
    dashed,
    draw=black!48
  },
  paperbox/.style={
    draw=black!58,
    line width=0.58pt,
    rounded corners=2pt,
    fill=white,
    align=center,
    inner xsep=5pt,
    inner ysep=5pt,
    minimum height=8mm
  },
  paperdata/.style={
    paperbox,
    fill=paperBlue,
    draw=paperBlueLine
  },
  paperproc/.style={
    paperbox,
    fill=paperTeal,
    draw=paperTealLine
  },
  papermodel/.style={
    paperbox,
    fill=paperAmber,
    draw=paperAmberLine,
    line width=0.65pt
  },
  paperout/.style={
    paperbox,
    fill=paperRose,
    draw=paperRoseLine,
    line width=0.65pt,
    font=\small\bfseries
  },
  paperstate/.style={
    paperbox,
    fill=paperViolet,
    draw=paperVioletLine
  },
  papernote/.style={
    font=\scriptsize,
    align=center,
    text=black!70,
    fill=white,
    inner sep=1.4pt
  },
  paperaxis/.style={
    line width=0.45pt,
    draw=black!72
  },
  papergrid/.style={
    line width=0.25pt,
    draw=black!15
  },
  paperplot/.style={
    line width=0.95pt,
    draw=paperBlueLine
  },
  paperplotA/.style={
    line width=0.95pt,
    draw=paperBlueLine
  },
  paperplotB/.style={
    line width=0.95pt,
    draw=paperAmberLine
  },
  papermark/.style={
    circle,
    draw=paperBlueLine,
    fill=white,
    line width=0.55pt,
    inner sep=1.6pt
  },
  papermarkA/.style={
    papermark,
    draw=paperBlueLine,
    fill=paperBlue
  },
  papermarkB/.style={
    papermark,
    draw=paperAmberLine,
    fill=paperAmber
  }
}

\resizebox{0.98\linewidth}{!}{%
\begin{tikzpicture}[paperfig, x=1mm, y=1mm]
  \node[paperdata, text width=16mm, minimum height=11mm] (task) at (7,38) {任务\\$x_i$};
  \draw[paperBlueLine, line width=0.45pt, fill=white] (1,42) rectangle (6,46);
  \draw[paperBlueLine, line width=0.45pt, fill=white] (3,40) rectangle (8,44);
  \draw[paperBlueLine, line width=0.45pt, fill=paperBlue] (5,38) rectangle (10,42);

  \node[paperstate, text width=17mm, minimum height=11mm] (budget) at (48,38) {预算\\档位 $b$};
  \draw[paperVioletLine, line width=0.6pt] (42,45) -- (54,45);
  \foreach \x in {42,46,50,54} \draw[paperVioletLine, line width=0.5pt] (\x,43.8) -- (\x,46.2);
  \fill[paperVioletLine] (50,45) circle (1.1);

  \node[papermodel, text width=28mm, minimum height=17mm] (orch) at (27.5,25) {编排器\\\scriptsize 按档位解锁动作};
  \draw[paperAmberLine, line width=0.5pt] (15.5,28.8) -- (39.5,28.8);
  \foreach \x/\lab in {19.5/生成,27.5/修补,35.5/扩展} {
    \draw[paperAmberLine, line width=0.45pt, fill=paperAmber!45, rounded corners=1pt]
      (\x-3,20.7) rectangle (\x+3,24.4);
    \node[font=\tiny] at (\x,22.6) {\lab};
  }

  \node[paperproc, text width=27mm, minimum height=15mm] (exec) at (27.5,7) {统一执行器 $E$\\\scriptsize 测试、超时、判定固定};
  \draw[paperTealLine, line width=0.5pt, rounded corners=1pt] (17,10.1) rectangle (38,13.8);
  \draw[paperTealLine, line width=0.5pt] (19,12) -- (22,12);
  \draw[paperTealLine, line width=0.5pt] (24,12) -- (31,12);
  \draw[paperTealLine, line width=0.5pt] (19,8.6) -- (35,8.6);

  \node[paperout, text width=15mm] (pass) at (6,-10) {判定\\$z_i^{(b)}$};
  \node[paperout, text width=15mm] (calls) at (27.5,-10) {调用\\$c_i^{(b)}$};
  \node[paperout, text width=15mm] (tokens) at (49,-10) {令牌\\$t_i^{(b)}$};

  \draw[paperline] (task) -- (orch);
  \draw[paperline] (budget) -- (orch);
  \draw[paperline] (orch) -- (exec);
  \draw[paperline] (exec) -- (pass);
  \draw[paperline] (exec) -- (calls);
  \draw[paperline] (exec) -- (tokens);

  \node[papernote, text width=44mm] at (27.5,-22)
    {同一评测口径下记录收益与成本，比较不同预算形态};
\end{tikzpicture}%
}
\caption{统一预算记账协议}
\label{fig:budget_accounting_protocol}
\end{figure}
